% Volume X: Normative Conclusions
% Part XVIII. Freedom Under Causal Opacity
% Chapter 108: Publicity of Power, Privacy of Persons
% Spine: proof-spines/ch108-spine.md

\chapter{Publicity of Power, Privacy of Persons}
\label{ch:ch108}

The normative orientation of this entire work can be stated simply.

\begin{quote}
\emph{Institutions exercising coercive power should be increasingly inspectable as their power grows. Persons should not become increasingly exposed merely because institutions acquire greater capacity to observe them.}
\end{quote}

The asymmetry is justified because institutions claim authority over others, whereas persons do not ordinarily owe the public exhaustive legibility. Publicity should therefore rise toward power, while privacy protects the ordinary conditions of personhood from becoming raw material for ever more capable systems of observation.
